\documentclass{article}


% 中文支持
\usepackage[UTF8,heading = true]{ctex}


% 页面布局
\usepackage[a4paper]{geometry}
\usepackage{fancyhdr}
\pagestyle{fancy}
\newcommand{\Lhead}[1]{\fancyhead[L]{#1}}
\newcommand{\Chead}[1]{\fancyhead[C]{#1}}
\newcommand{\Rhead}[1]{\fancyhead[R]{#1}}
\newcommand{\Lfoot}[1]{\fancyfoot[L]{#1}}
\newcommand{\Cfoot}[1]{\fancyfoot[C]{#1}}
\newcommand{\Rfoot}[1]{\fancyfoot[R]{#1}}
\renewcommand{\headrulewidth}{0.4pt}
\renewcommand{\footrulewidth}{0pt}
\usepackage{lastpage}
%待填信息
% -----------------------------------------------------
\newcommand{\Course}{深度学习平台与应用}             		%      |
\newcommand{\Assignment}{深度学习作业 5}		  %      | 
\newcommand{\Name}{田永铭}				    	 %      |
\newcommand{\ID}{221900180}				  	   %       |
\newcommand{\Email}{tianym2022@smail.nju.edu.cn}  %    |
% ------------------------------------------------------
\fancyhead[L]{\Course}
\fancyhead[C]{\Assignment}
\fancyhead[R]{\ID \Name}
\fancyfoot[C]{\thepage\ / \pageref*{LastPage}}	
\setlength\headheight{36pt}
% 自定义页面样式第一页只有页脚没有页眉
\makeatletter
\def\ps@onlyfoot{
	\let\@oddhead\@empty
	\let\@evenhead\@empty
	\def\@oddfoot{\hfil\thepage\ / \pageref*{LastPage}\hfil}
	\def\@evenfoot{\hfil\thepage\ / \pageref*{LastPage}\hfil}
}
\makeatother
% maketitle
\title{\Huge\Assignment}
\author{\ID \Name \quad\href{mailto:tianym2022@smail.nju.edu.cn}{\color{magenta}\Email}}
\date{\today}


% 字体和颜色
\usepackage{fontspec} 
\fontsize{12pt}{\baselineskip}\selectfont
\usepackage{xcolor}


% 数学支持
\usepackage{amssymb}
\usepackage{amsfonts}
\usepackage{amsmath}
\usepackage{esint}
\usepackage{gensymb}
\usepackage{mathtools}
\usepackage{amsthm}

% 左边大括号
%\[
%\left\{
%\begin{aligned}
%	& \text{approximate  covariance matrix of GP} \\
%	& \text{only use subset}\\
%	& \text{conditional independence assumption}
%\end{aligned}
%\right.
%\]

% 矩阵
%\[
%\begin{matrix}
%	a & b \\
%	c & d
%\end{matrix}
%\]

% 行列式
%\[
%\begin{vmatrix}
%	a & b \\
%	c & d
%\end{vmatrix}
%\]


% 标题格式
\usepackage{titlesec}
% 自定义 section 标题格式
\renewcommand{\thesection}{\chinese{section}、}
\titleformat{\section}
{\normalfont\Large\bfseries\centering\color[rgb]{0,0,0.8}} % 设定字体和大小
{\thesection}{0em} % 设定编号和标题之间的距离
{}
% 自定义 subsection 标题格式
\renewcommand{\thesubsection}{\arabic{section}.\arabic{subsection} }
\titleformat{\subsection}
{\normalfont\large\bfseries\color[rgb]{0.23,0.65,0.99}}
{\thesubsection}{0em}
{}
% 自定义 subsubsection 标题格式
\renewcommand{\thesubsubsection}{\arabic{section}.\arabic{subsection}.\arabic{subsubsection} }
\titleformat{\subsubsection}
{\normalfont\normalsize\bfseries}
{\thesubsubsection}{0em}
{}
% 设置 section 标题的间距
\titlespacing*{\section}
{0pt}{3.5ex plus 1ex minus .2ex}{2.3ex plus .2ex}
% 设置 subsection 标题的间距
\titlespacing*{\subsection}
{0pt}{3.25ex plus 1ex minus .2ex}{1.5ex plus .2ex}
% 设置 subsubsection 标题的间距
\titlespacing*{\subsubsection}
{0pt}{3.25ex plus 1ex minus .2ex}{1.5ex plus .2ex}


% 代码 
% 如非默认，需使用[language=xxx,style=mystyle]
\usepackage{listings}
% language:  {[x86masm]Assembler}(汇编); python; C; C++; Java; SQL; Matlab; R; bash; make(makefile)
\lstdefinestyle{mystyle}{
	basicstyle=\ttfamily\small,
	numbers=left,
	numberstyle=\tiny\color{magenta},
	stepnumber=1,
	numbersep=10pt,
	backgroundcolor=\color[rgb]{0.95,0.95,0.95},
	showspaces=false,
	showstringspaces=false,
	showtabs=false,
	frame=single,
	rulecolor=\color{black},
	tabsize=2,
	captionpos=b,
	breaklines=true,
	breakatwhitespace=false,
	title=\lstname,
	keywordstyle=\color{blue},
	commentstyle=\color[rgb]{0.0,0.5,0.0},
	stringstyle=\color[rgb]{0.23,0.65,0.99}\itshape,
	escapeinside={\%*}{*)},
	morekeywords={}, % 添加更多关键字
	literate=
	{[}{{{\color{red}[}}}1
	{]}{{{\color{red}]}}}1
	{\{}{{{\color[rgb]{0.9,0,0.9}\{}}}1
	{\}}{{{\color[rgb]{0.9,0,0.9}\}}}}1
%	{0}{{{\color[rgb]{0.2,0.9,0.4}0}}}1
%	{1}{{{\color[rgb]{0.2,0.9,0.4}1}}}1
%	{2}{{{\color[rgb]{0.2,0.9,0.4}2}}}1
%	{3}{{{\color[rgb]{0.2,0.9,0.4}3}}}1
%	{4}{{{\color[rgb]{0.2,0.9,0.4}4}}}1
%	{5}{{{\color[rgb]{0.2,0.9,0.4}5}}}1
%	{6}{{{\color[rgb]{0.2,0.9,0.4}6}}}1
%	{7}{{{\color[rgb]{0.2,0.9,0.4}7}}}1
%	{8}{{{\color[rgb]{0.2,0.9,0.4}8}}}1
%	{9}{{{\color[rgb]{0.2,0.9,0.4}9}}}1
}
\lstset{style=mystyle}


% 图表
\usepackage{graphicx}
\usepackage{booktabs}
\usepackage{caption}
\captionsetup{font={small}}
\usepackage{diagbox}
\usepackage{array}
\usepackage{pgfplots}
\pgfplotsset{compat=1.17}

% 表1
%\begin{table}[h!]
%	\centering
%	\caption{Performance comparison on mini-ImageNet and tiered-ImageNet datasets.}
%	\begin{tabular}{@{}lcccccccc@{}}
%		\toprule
%		Dataset & Method & Backbone & 5-way 1-shot \\
%		\midrule
%		\multirow{9}{*}{\rotatebox[origin=c]{90}{mini-ImageNet}} 
%		& ProtoNet & Conv-4 & 46.291 \\
%		& MAML & Conv-4 & 49.704 \\
%		& BOIL & Conv-4 & 48.456 \\
%		& CAN & ResNet-12 & 59.493\\
%		& LEO & ResNet-12 & 55.056 \\
%		& SKD & ResNet-12 & 66.700  \\
%		& R2D2 & ResNet-12 & 59.802 \\
%		& DiffKendall (Ours) & ResNet-12 & \textbf{68.727} \\
%		\midrule
%		\multirow{6}{*}{\rotatebox[origin=c]{90}{tiered-ImageNet}} 
%		& ProtoNet & Conv-4 & 46.653 \\
%		& Relation Networks & Conv-4 & 55.107  \\
%		& MAML & Conv-4 & 51.038 \\
%		& R2D2 & ResNet-12 & 64.916  \\
%		& DiffKendall (Ours) & ResNet-12 & \textbf{66.598} \\
%		\bottomrule
%	\end{tabular}
%\end{table}

% 表2
%\begin{table}[t]
%	\caption{Classification accuracies for naive Bayes and flexible
%		Bayes on various data sets.}
%	\label{sample-table}
%	\vskip 0.15in
%	\begin{center}
%		\begin{small}
%			\begin{sc}
%				\begin{tabular}{lcccr}
%					\toprule
%					Data set & Naive & Flexible & Better? \\
%					\midrule
%					Breast    & 95.9$\pm$ 0.2& 96.7$\pm$ 0.2& $\surd$ \\
%					Cleveland & 83.3$\pm$ 0.6& 80.0$\pm$ 0.6& $\times$\\
%					Glass2    & 61.9$\pm$ 1.4& 83.8$\pm$ 0.7& $\surd$ \\
%					Credit    & 74.8$\pm$ 0.5& 78.3$\pm$ 0.6&         \\
%					Horse     & 73.3$\pm$ 0.9& 69.7$\pm$ 1.0& $\times$\\
%					Meta      & 67.1$\pm$ 0.6& 76.5$\pm$ 0.5& $\surd$ \\
%					Pima      & 75.1$\pm$ 0.6& 73.9$\pm$ 0.5&         \\
%					Vehicle   & 44.9$\pm$ 0.6& 61.5$\pm$ 0.4& $\surd$ \\
%					\bottomrule
%				\end{tabular}
%			\end{sc}
%		\end{small}
%	\end{center}
%	\vskip -0.1in
%\end{table}


% 算法
%\usepackage{algorithm}
%\usepackage{algorithmic}
%\begin{algorithm}[tb]
%	\caption{Bubble Sort}
%	\label{alg:example}
%	\begin{algorithmic}
	%		\STATE {\bfseries Input:} data $x_i$, size $m$
	%		\REPEAT
	%		\STATE Initialize $noChange = true$.
	%		\FOR{$i=1$ {\bfseries to} $m-1$}
	%		\IF{$x_i > x_{i+1}$}
	%		\STATE Swap $x_i$ and $x_{i+1}$
	%		\STATE $noChange = false$
	%		\ENDIF
	%		\ENDFOR
	%		\UNTIL{$noChange$ is $true$}
	%	\end{algorithmic}
%\end{algorithm}

% 超链接
\usepackage{hyperref}


% 论文参考文献
% 需在正文结尾加入 \bibliography{references}，在同目录下创建references.bib
\bibliographystyle{plain}


% 表情
% https://ctan.math.washington.edu/tex-archive/graphics/pgf/contrib/tikzsymbols/tikzsymbols.pdf
\usepackage{tikzsymbols}
\usepackage{tikz}
% 背景颜色：\colorbox{yellow}{...}
% 3D版本： 在前面加上d,比如：dSmiley

% 常用
% 笑脸：\Smiley[⟨scale⟩][⟨color⟩]
% 哭脸：\Sadey[⟨scale⟩][⟨color⟩]
% 生气脸： \Annoey[⟨scale⟩][⟨color⟩]
% 大笑： \Laughey[⟨scale⟩][⟨color⟩][⟨mouth color⟩]
% 挤眉弄眼： \Winkey[⟨scale⟩][⟨color⟩]
% 去世： \Xey[⟨scale⟩][⟨color⟩]
% 升天： \wInnocey[⟨scale⟩]
% 口吐芬芳： \Vomey[⟨scale⟩][⟨color⟩][⟨vomit color⟩]
% 猫： \Cat[⟨scale⟩]
% 炸弹： \Ninja[⟨scale⟩][⟨color⟩][⟨headband color⟩][⟨eye color⟩]
% 困： \Sleepey[⟨scale⟩][⟨color⟩][⟨cap color⟩][⟨star color⟩]
% 口罩：\Maskey[⟨scale⟩][⟨color⟩][⟨mask color⟩]
% 渐变弧度脸(-2到2)： \Changey[⟨scale⟩][⟨color⟩]{⟨mood⟩}
% 树（春夏秋冬）： \Springtree[⟨scale⟩]
% 咖啡： \Coffeecup[⟨scale⟩]



% 简化命令
\newcommand{\hs}[1]{\hspace{{#1}em}} %空格
\newcommand{\qd}{\quad} %空格
\newcommand{\qqd}{\qquad} %空格
% 链接
\newcommand{\myhref}[2]{\href{#1}{{\color[rgb]{0.9,0,0.9}{#2}}}}
\newcommand{\myred}[1]{{\color{red}{#1}}} %红色
% 图 args：figures下图片路径，大小，图片显示名
\newcommand{\myfig}[3]{
	\begin{figure}[h!]
		\centering
		\includegraphics[width=#2\linewidth]{figures/#1}
		\caption{#3}
		\label{fig:#1}
	\end{figure}
}
% 左右分栏插入图片
% args: 左图文件名，左图宽度，左图标题，右图文件名，右图宽度，右图标题
\newcommand{\mytwocolumnfig}[6]{
	\begin{figure}[h!]
		\centering
		\begin{minipage}[h!]{0.45\linewidth}
			\centering
			\includegraphics[width=#2\linewidth]{figures/#1}
			\caption{#3}
			\label{fig:#1}
		\end{minipage}
		\hfill
		\begin{minipage}[h!]{0.45\linewidth}
			\centering
			\includegraphics[width=#5\linewidth]{figures/#4}
			\caption{#6}
			\label{fig:#4}
		\end{minipage}
	\end{figure}
}


\newtheoremstyle{Problem}
{}
{}
{}
{}
{\bfseries}
{}
{ }
{Problem\thmnumber{ #2}\thmname{ #1}\thmnote{ (#3)} }
\theoremstyle{Problem}

\newtheorem{prob}{}
\newtheoremstyle{Solution}
{}
{}
{}
{}
{\bfseries}
{:}
{ }
{\thmname{#1}}
\makeatletter
\makeatother
\theoremstyle{Solution}
\newtheorem*{sol}{Solution}


\begin{document}
	\maketitle
	
	\begin{prob}
		\textbf{列举几种常见的激活函数，并简要说明它们的作用。}
	\end{prob}
	
	\begin{sol}
	解：
\begin{enumerate}
	\item \textbf{Sigmoid 函数} \\
	\textit{公式：} \( \sigma(x) = \frac{1}{1 + e^{-x}} \) \\
	\textit{作用：} 将输入值映射到 0 和 1 之间。适用于二分类问题的输出层，因为其输出可解释为概率。但在深层网络中容易导致梯度消失问题，因此不常用于隐藏层。
	
	\item \textbf{Tanh（双曲正切）函数} \\
	\textit{公式：} \( \tanh(x) = \frac{e^{x} - e^{-x}}{e^{x} + e^{-x}} \) \\
	\textit{作用：} 输出值在 -1 和 1 之间，比 Sigmoid 有更好的性能，因为其输出是零中心的，有助于加速梯度下降。与 Sigmoid 类似，Tanh 也存在梯度消失的问题。
	
	\item \textbf{ReLU（Rectified Linear Unit）} \\
	\textit{公式：} \( \text{ReLU}(x) = \max(0, x) \) \\
	\textit{作用：} 当输入大于 0 时，输出为输入值；否则输出为 0。ReLU 是目前最常用的激活函数之一，能够有效缓解梯度消失问题，且计算简单。其缺点是会出现“死神经元”问题，导致部分神经元永远不会激活。
	
	\item \textbf{Leaky ReLU} \\
	\textit{公式：} \( \text{Leaky ReLU}(x) = \max(\alpha x, x) \)（其中 \(\alpha\) 是一个小常数） \\
	\textit{作用：} Leaky ReLU 是对 ReLU 的改进，允许输入小于 0 时输出为一个很小的负数，避免了死神经元问题。通常，\(\alpha\) 取值较小，如 0.01。
	
	\item \textbf{GELU (Gaussian Error Linear Unit)} \\
	\textit{公式：} \( \text{GELU}(x) = 0.5 x \left( 1 + \tanh\left( \sqrt{\frac{2}{\pi}} (x + 0.044715 x^3) \right) \right) \) \\
	\textit{作用：} GELU 是一种平滑的非线性激活函数，常用于 Transformer 等模型中。它结合了高斯误差函数和线性激活的优点，使得网络在训练时能够获得更好的梯度流，避免 ReLU 的死神经元问题，并具有类似于 ReLU 的高效性。
\end{enumerate}
	\end{sol}
	
	\newpage
	
	\begin{prob}
		\textbf{什么是反向传播？它的核心思想是什么？}
	\end{prob}
	
	\begin{sol}
		解:\\
		\textbf{反向传播（Backpropagation） :}\\是一种用于训练神经网络的算法，它通过链式法则有效地计算损失函数关于网络中每个参数的梯度，从而实现模型的优化。反向传播通过梯度下降算法来更新神经网络中的权重和偏置，使得损失函数的值逐渐减小，从而提高网络的性能。\\
		\textbf{核心思想：}\\
		反向传播的核心思想是利用\textbf{链式法则}，从输出层开始，\textbf{逐层将误差反向传播到隐藏层和输入层}，计算每个参数对损失函数的\textbf{偏导数}，并利用这些梯度更新参数。
	\end{sol}
	
	\begin{prob}
		\textbf{Numpy 实现反向传播梯度。}
		\begin{lstlisting}[language=python,style=mystyle]
import numpy as np

# 设置随机种子，方便复现
np.random.seed(42)

# 初始化参数
W1 = np.random.randn(2, 2)  # 输入层到隐藏层的权重
b1 = np.random.randn(2)  # 隐藏层的偏置
W2 = np.random.randn(2)  # 隐藏层到输出层的权重
b2 = np.random.randn()  # 输出层的偏置

# 输入数据和真实标签
x = np.array([1.0, -1.5])
y_true = 0.5

# 学习率
learning_rate = 0.01


# ReLU
def relu(x):
	#todo 输入x 返回relu(x)
	pass


# ReLU 的梯度
def relu_derivative(x):
	return np.where(x > 0, 1, 0) 


# 均方损失函数和其导数
def mse_loss(y_pred, y_true):
	#todo 计算均方误差
	pass


def mse_loss_derivative(y_pred, y_true):
	return y_pred - y_true  # 对 y_pred 求导


# 训练过程
epochs = 100
for epoch in range(epochs):
	# 前向传播
	
	z1 = np.dot(W1, x) + b1  
	h = relu(z1) 
	y_pred = np.dot(W2, h) + b2 
	
	# 计算损失
	loss = mse_loss(y_pred, y_true)
	
	# 反向传播
	d_loss_y_pred = mse_loss_derivative(y_pred, y_true)  # 损失对预测值的梯度
	#todo 
	#d_loss_W2 = ?  对 W2 的梯度
	#d_loss_b2 = ?  对 b2 的梯度
	
	#d_loss_h =?  损失对隐藏层输出 h 的梯度
	#d_loss_z1 =?  隐藏层的梯度 (通过 ReLU 计算)
	
	#d_loss_W1 =?  对 W1 的梯度
	#d_loss_b1 =?  对 b1 的梯度
	
	# 参数更新 (随机梯度下降)
	W1 -= learning_rate * d_loss_W1
	b1 -= learning_rate * d_loss_b1
	W2 -= learning_rate * d_loss_W2
	b2 -= learning_rate * d_loss_b2
	
	# 打印训练过程中的损失
	if epoch % 10 == 0:
	print(f"Epoch {epoch}: Loss = {loss}")

# 打印最终的损失值
print(f"Final Loss: {loss}")
		\end{lstlisting}
	\end{prob}
	
	\begin{sol}
		补充部分代码如下：
		\begin{lstlisting}[language=python,style=mystyle]
		
# ReLU
def relu(x):
	#todo 输入x 返回relu(x)
	return np.maximum(0, x)		
......
		
def mse_loss(y_pred, y_true):
	#todo 计算均方误差
	return np.mean((y_pred - y_true) ** 2) / 2		
......

# 反向传播
d_loss_y_pred = mse_loss_derivative(y_pred, y_true)  # 损失对预测值的梯度
# d_loss_W2: 损失对 W2 的梯度
d_loss_W2 = d_loss_y_pred * h.T   # 对 W2 求导，链式法则，h 是列向量，使用 h.T 进行矩阵乘法
# d_loss_b2: 损失对 b2 的梯度
d_loss_b2 = d_loss_y_pred  # 对 b2 求导
# d_loss_h: 损失对隐藏层输出 h 的梯度
d_loss_h =  d_loss_y_pred * W2.T # 反向传播到隐藏层
# d_loss_z1: 隐藏层的梯度 (通过 ReLU 计算)
d_loss_z1 = d_loss_h * relu_derivative(z1)  # 通过 ReLU 的导数来计算
# d_loss_W1: 损失对 W1 的梯度
d_loss_W1 = np.outer(d_loss_z1, x)  # 对 W1 求导，链式法则，x 是输入
# d_loss_b1: 损失对 b1 的梯度
d_loss_b1 = d_loss_z1  # 对 b1 求导
		\end{lstlisting}
		
	其它求导步骤都很简单，利用链式法则即可，对W1的梯度，由于W1是一个矩阵，所以求导比较复杂，代码24行具体解释如图~\ref{fig:2}。（注意这里应该每一步都有转置和每一步都没有出来是一样的，本质相同。）
	
	\begin{figure}[h!]
		\centering
		\includegraphics[width=1\linewidth]{Figures/2}
		\caption{}
		\label{fig:2}
	\end{figure}
	
	以上是全写成标量计算的，其实也能用正常方法，$x$在这里是数组形状$[x_1,x_2]$，$dz$也是$[dz_1,dz_2]$，只要计算$x^\top dz$(矩阵乘法)，也能得到一样的结果。
	
		
	代码运行结果如图~\ref{fig:1}, 可以看到损失逐渐降低，最终很小，说明训练有效。
		
\begin{figure}[h!]
	\centering
	
\includegraphics[width=0.5\linewidth]{Figures/1}
	\caption{tarin\_loss}
	\label{fig:1}
\end{figure}
	\end{sol}
	
	
	\begin{prob}
		\textbf{什么是全连接网络（FC）或多层感知机（MLP）？它与线性分类器的主要区别是什么？}
	\end{prob}
	
	\begin{sol}
		解：
		
		
		\textbf{全连接网络（FC）或多层感知机（MLP）}
		
		全连接网络（FC）或多层感知机（MLP）是一种神经网络模型，由多层神经元组成，其中每一层的每个神经元与前一层的每个神经元都有连接。它是最基础的神经网络类型之一，广泛用于分类和回归任务。
		\begin{itemize}
			\item \textbf{输入层}：接收原始输入数据。
			\item \textbf{隐藏层}：通过加权和偏置进行计算，并通过激活函数进行非线性变换。
			\item \textbf{输出层}：根据任务需求输出结果，如分类标签或回归值。
		\end{itemize}
		
		MLP 是一种前馈神经网络，信息在网络中单向传播，不会有反馈连接。每一层的输出都成为下一层的输入。
		
		\textbf{线性分类器}
		
		线性分类器（如线性回归或支持向量机）通过输入特征的线性组合来做出预测，输出一个线性函数，通常表示为：
		\[
		y = w_1 x_1 + w_2 x_2 + \cdots + w_n x_n + b
		\]
		其中 \( x_1, x_2, \dots, x_n \) 是输入特征，\( w_1, w_2, \dots, w_n \) 是对应的权重，\( b \) 是偏置项。
		
		而多层感知机通过多个层级和非线性激活函数的组合，可以拟合更复杂的非线性关系。每一层的输出通过激活函数进行非线性变换，从而具有更强的表达能力。
		
		\textbf{主要区别}
		\begin{itemize}
			\item \textbf{线性与非线性}：
			\begin{itemize}
				\item 线性分类器只能表示输入特征的线性关系。
				\item MLP 使用非线性激活函数，使其能够学习复杂的非线性关系。
			\end{itemize}
			\item \textbf{层次结构}：
			\begin{itemize}
				\item 线性分类器只有一个层次，直接从输入到输出。
				\item MLP 由多层组成，包括至少一个隐藏层。
			\end{itemize}
			\item \textbf{表达能力}：
			\begin{itemize}
				\item 线性分类器的表达能力较弱，一般仅适用于线性可分的任务。
				\item MLP 可以逼近任何复杂的函数，适用于更广泛的任务。
			\end{itemize}
		\end{itemize}
	\end{sol}
	
	\begin{prob}
		\textbf{在梯度流的基本计算模式中，加法门(add gate)、乘法门(mul gate)和最大值门(max gate)分别具有什么特点？}
	\end{prob}
	
	\begin{sol}
		解：\\
		\textbf{加法门：}加法门不改变梯度，只是将梯度传递到所有输入。
		
		\textbf{乘法门：}乘法门会将梯度乘以其他输入的乘积（假设是多个，如果是两个输入的乘法则是另一个输入）传到这个输入。
		
		\textbf{最大值门：}最大值门会将梯度传递给选择出的最大值，但如果最大值有多个，则平分梯度。
		
		直观展示见图~\ref{fig:3}。
		
\begin{figure}[h!]
	\centering
	\includegraphics[width=1\linewidth]{Figures/3}
	\caption{}
	\label{fig:3}
\end{figure}
		
		
	\end{sol}
	

\end{document}
